% ============================================================================
%  math-environments.tex
%  Ambienti matematici con amsthm.
%
%  Versione adeguata alle richieste redazionali (Springer):
%   - corpo del testo a dimensione NORMALE (niente \footnotesize);
%   - nessun rientro a destra (margine sinistro azzerato);
%   - "quadratino" di fine sezione, giustificato a destra, per delimitare
%     in modo univoco Esempio / Esercizio / Osservazione (l'ambiente proof
%     usa gia` il quadratino tramite \qed di amsthm).
%
%  >>> PER ADEGUARE ALLO STILE TIPOGRAFICO: agire solo sul blocco "LEVE DI
%      STILE" qui sotto. Il resto del file non va modificato. <<<
% ============================================================================

% ----------------------------------------------------------------------------
%  LEVE DI STILE  (unico punto da toccare per il template editoriale)
% ----------------------------------------------------------------------------
% Dimensione del corpo degli ambienti evidenziati.
%   Default: \normalsize (richiesta Springer). Per tornare al corpo ridotto
%   basta porre \renewcommand{\ambientebodysize}{\footnotesize}.
\newcommand{\ambientebodysize}{\normalsize}

% Rientro sinistro degli ambienti evidenziati.
%   Default: 0pt (nessuno spostamento a destra, richiesta Springer).
\newlength{\ambienteleftmargin}
\setlength{\ambienteleftmargin}{0pt}

% Quadratini di fine ambiente (vuoto = \square, pieno = \blacksquare).
%   Convenzione adottata: pieno per Esempio/Esercizio, vuoto per Osservazione.
\newcommand{\finemarkesempio}{\ensuremath{\blacksquare}}
\newcommand{\finemarkesercizio}{\ensuremath{\blacksquare}}
\newcommand{\finemarkoss}{\ensuremath{\square}}
% ----------------------------------------------------------------------------


% ESEMPIO ---------------------------------------------------------------------
\newtheoremstyle{esempiostyle}
{6pt}{6pt}
{\ambientebodysize}
{0pt}
{\bfseries}
{.}
{ }
{\bfseries #1~#2\ifx#3\empty\else~(#3)\fi}
\theoremstyle{esempiostyle}
\newtheorem{esempioinner}{Esempio}[chapter]
\newenvironment{esempio}[1][]{%
	\begin{list}{}{%
			\setlength{\leftmargin}{\ambienteleftmargin}%
			\setlength{\rightmargin}{0pt}%
			\setlength{\labelwidth}{0pt}%
			\setlength{\labelsep}{0pt}%
			\setlength{\itemindent}{0pt}%
			\setlength{\topsep}{4pt}%
			\setlength{\parsep}{0pt}}%
		\item
		\renewcommand{\qedsymbol}{\finemarkesempio}%
		\begin{esempioinner}[#1]%
			\pushQED{\qed}\ignorespaces
		}{%
			\popQED
		\end{esempioinner}%
	\end{list}%
}

% ESERCIZIO -------------------------------------------------------------------
\newtheoremstyle{esercziostyle}
{6pt}{6pt}
{\ambientebodysize}
{0pt}
{\bfseries}
{.}
{ }
{\bfseries #1~#2\ifx#3\empty\else~(#3)\fi}
\theoremstyle{esercziostyle}
\newtheorem{esercizioinner}{Esercizio}[chapter]
\newenvironment{esercizio}[1][]{%
	\begin{list}{}{%
			\setlength{\leftmargin}{\ambienteleftmargin}%
			\setlength{\rightmargin}{0pt}%
			\setlength{\labelwidth}{0pt}%
			\setlength{\labelsep}{0pt}%
			\setlength{\itemindent}{0pt}%
			\setlength{\topsep}{4pt}%
			\setlength{\parsep}{0pt}}%
		\item
		\renewcommand{\qedsymbol}{\finemarkesercizio}%
		\begin{esercizioinner}[#1]%
			\pushQED{\qed}\ignorespaces
		}{%
			\popQED
		\end{esercizioinner}%
	\end{list}%
}


% SVOLGIMENTO -----------------------------------------------------------------
% Corpo normale (era \footnotesize). Nessun quadratino: lo svolgimento e`
% interno all'esempio/esercizio, il cui quadratino chiude l'intero blocco.
\newenvironment{svolgimento}{%
	\begin{list}{}{%
			\setlength{\leftmargin}{0cm}%
			\setlength{\rightmargin}{0cm}%
			\setlength{\topsep}{2pt}%
			\setlength{\parsep}{0pt}}%
		\item\textit{Svolgimento.}\space\ambientebodysize%
	}{%
	\end{list}%
}

%--------------------------------------------------------------------
% \esottotit{...} — sotto-titolo interno agli esempi
% Produce un sotto-titolo in grassetto a riga propria, con piccolo
% spazio sopra per separare dal blocco precedente.
% Uso:  \esottotit{Problema virtuale.} Per isolare la grandezza ...
%--------------------------------------------------------------------
\newcommand{\esottotit}[1]{\par\smallskip\noindent\textbf{#1}}

%--------------------------------------------------------------------
% \risultato — etichetta del risultato negli esercizi
%--------------------------------------------------------------------
\newcommand{\risultato}{\par\smallskip\noindent\textit{Risultato.}\space}


% ABSTRACT (gia` "Sommario") --------------------------------------------------
% Richiesta Springer: dicitura inglese "Abstract", testo in italiano.
% Corpo normale e nessun rientro (lo stile dell'abstract sara` comunque
% definito dal template editoriale Springer).
\newenvironment{sommario}{%
	\begin{list}{}{%
			\setlength{\leftmargin}{0cm}%
			\setlength{\rightmargin}{0cm}%
			\setlength{\topsep}{4pt}%
			\setlength{\parsep}{0pt}}%
		\item\bfseries Abstract.\normalfont\itshape\space
	}{%
	\end{list}%
}



% OSSERVAZIONE ----------------------------------------------------------------
\newtheoremstyle{ossstyle}
{6pt}{6pt}
{\upshape\ambientebodysize}
{0pt}
{\bfseries}
{.}
{ }
{\bfseries #1~#2\ifx#3\empty\else~(#3)\fi}
\theoremstyle{ossstyle}
\newtheorem{ossinner}{Osservazione}[chapter]
\newenvironment{oss}[1][]{%
	\begin{list}{}{%
			\setlength{\leftmargin}{\ambienteleftmargin}%
			\setlength{\rightmargin}{0pt}%
			\setlength{\labelwidth}{0pt}%
			\setlength{\labelsep}{0pt}%
			\setlength{\itemindent}{0pt}%
			\setlength{\topsep}{4pt}%
			\setlength{\parsep}{0pt}}%
		\item
		\renewcommand{\qedsymbol}{\finemarkoss}%
		\begin{ossinner}[#1]%
			\pushQED{\qed}\ignorespaces
		}{%
			\popQED
		\end{ossinner}%
	\end{list}%
}

% ── Teorema, Lemma, Proposizione, Corollario ─────────────────────────────
\newtheoremstyle{teoremastyle}
{6pt}{6pt}
{\itshape\ambientebodysize}
{0pt}
{\bfseries}
{.}
{ }
{\bfseries #1~#2\ifx#3\empty\else~(#3)\fi}
\theoremstyle{teoremastyle}
\newtheorem{teorema}{Teorema}[section]
\newtheorem{lemma}[teorema]{Lemma}
\newtheorem{proposizione}[teorema]{Proposizione}
\newtheorem{corollario}[teorema]{Corollario}

% ── Definizione ──────────────────────────────────────────────────────────
\newtheoremstyle{definizionestyle}
{6pt}{6pt}
{\normalsize}
{0pt}
{\bfseries}
{.}
{ }
{\bfseries #1~#2\ifx#3\empty\else~(#3)\fi}
\theoremstyle{definizionestyle}
\newtheorem{definizione}[teorema]{Definizione}

% ── Osservazione (numerata con teorema) ──────────────────────────────────
\newtheoremstyle{osservazionestyle}
{6pt}{6pt}
{\normalsize}
{0pt}
{\bfseries}
{.}
{ }
{\bfseries #1~#2\ifx#3\empty\else~(#3)\fi}
\theoremstyle{osservazionestyle}
\newtheorem{osservazione}[teorema]{Osservazione}
\newtheorem*{nota}{Nota}

% ── Ambiente proof ───────────────────────────────────────────────────────
% Corpo normale e nessun rientro a destra; conserva il quadratino \qed.
\makeatletter
\renewenvironment{proof}[1][\proofname]{%
	\begin{list}{}{%
			\setlength{\leftmargin}{\ambienteleftmargin}%
			\setlength{\rightmargin}{0pt}%
			\setlength{\topsep}{8pt}%
			\setlength{\parsep}{0pt}}%
		\item\par\ambientebodysize
		\pushQED{\qed}%
		\normalfont\ambientebodysize
		\topsep8pt \partopsep8pt
		\trivlist
		\item[\hskip\labelsep\itshape #1\@addpunct{.}]\ignorespaces
	}{%
		\popQED\endtrivlist\@endpefalse
	\end{list}%
}
\makeatother